\chapter*{I}
\label{chap:i}

\term{Ideal (Ring Theory)}
An ideal $I$ of a ring $R$ is an additive subgroup closed under multiplication by elements of $R$. A left ideal requires $ra\in I$, a right ideal requires $ar\in I$, and a two-sided ideal satisfies both. The quotient $R/I$ has a ring structure exactly when $I$ is two-sided. \par\noindent\textbf{Applications:} ideals describe divisibility, kernels of homomorphisms, and solution sets of polynomial equations.

\term{Ideal Class Group}
The ideal class group is the group of fractional ideals of a number field modulo principal fractional ideals. It measures the failure of unique factorization of elements: a trivial class group implies unique factorization of ideals and, in important cases, of elements. Its finite order is the class number. \par\noindent\textbf{Applications:} class groups organize arithmetic in number fields and appear in class field theory, Diophantine equations, and cryptography.

\term{Ideal Sheaf}
The sheaf $\mathcal{I}_Y$ on a space $X$ associated with a closed subscheme $Y$, whose sections over an open set $U$ are the functions in $\mathcal{O}_X(U)$ that vanish on $Y \cap U$. It is the kernel of the restriction map $\mathcal{O}_X \to i_* \mathcal{O}_Y$.

\term{Idempotent Element}
An element $e$ is idempotent if $e^2=e$. Besides $0$ and $1$, a ring may have idempotents that split its structure: for a central idempotent, $R\cong eR\times(1-e)R$. An idempotent matrix is a projection. \par\noindent\textbf{Applications:} idempotents decompose rings and modules and represent repeated decisions or projections in algebra and computation.

\term{Identity Element}
An element $e$ in a set $S$ with a binary operation $\cdot$ such that $e \cdot a = a \cdot e = a$ for all $a \in S$. In group theory, the identity is unique. In category theory, an identity morphism $\text{id}_X: X \to X$ exists for every object.

\term{Identity Theorem}
If two holomorphic functions on a connected domain agree on a set having a limit point inside the domain, they agree everywhere. Thus a holomorphic function is determined by surprisingly sparse data, and analytic continuation is unique. \par\noindent\textbf{Applications:} the theorem proves uniqueness in complex analysis and helps extend solutions of analytic differential equations.

\term{Idempotent Matrix}
A square matrix $P$ is idempotent when $P^2=P$. It projects vectors onto its image along its kernel; its only possible eigenvalues are $0$ and $1$, so it is diagonalizable. \par\noindent\textbf{Applications:} projection matrices are used in least squares, statistics, numerical linear algebra, and decomposing vector spaces.

\term{Iff}
Abbreviation for "if and only if", denoting logical equivalence between two statements.

\term{Image}
The image of $A$ under $f$ is $f(A)=\{f(x):x\in A\}$, the set of values actually attained. It is different from the whole codomain when $f$ is not onto. \par\noindent\textbf{Applications:} images describe reachable states, solution sets, and the range of linear transformations.

\term{Imaginary Number}
A purely imaginary number has the form $bi$ with $b\in\mathbb R$ and $i^2=-1$. It lies on the imaginary axis of the complex plane. Together with real numbers, imaginary numbers form the complex numbers $a+bi$. \par\noindent\textbf{Applications:} complex numbers solve polynomial equations and represent oscillations, AC circuits, waves, and Fourier modes.

\term{Imaginary Unit}
The number $i$ defined by $i^2 = -1$. It allows for the construction of the field of complex numbers $\mathbb{C}$. In physics, $i$ often appears in the Schrödinger equation and Fourier transforms.

\term{Immersion}
A smooth map $f:M\to N$ is an immersion if its differential $df_p$ is injective at every point. Locally it preserves tangent directions, but its image may cross itself; an embedding additionally requires suitable global one-to-one behavior. \par\noindent\textbf{Applications:} immersions describe curves and surfaces in differential geometry, computer graphics, and the kinematics of constrained systems.

\term{Implication}
A logical connective $P \Rightarrow Q$ true except when $P$ is true and $Q$ is false.

\term{Improper Fraction}
A fraction whose numerator is greater than or equal to its denominator.

\term{Improper Integral}
An improper integral has an unbounded interval or an unbounded integrand. For example, $\int_a^\infty f(x)\,dx=\lim_{b\to\infty}\int_a^b f(x)\,dx$ when the limit exists. Convergence must be checked rather than assumed. \par\noindent\textbf{Applications:} improper integrals measure total probability, energy, mass, and long-time effects in applied models.

\term{Implicit Function Theorem}
If $F:\mathbb R^{n+m}\to\mathbb R^m$ is smooth, $F(x_0,y_0)=0$, and the Jacobian $\partial F/\partial y$ is invertible there, then locally the equation $F(x,y)=0$ defines a unique smooth function $y=g(x)$. It explains when dependent variables can be solved from constraints. \par\noindent\textbf{Applications:} it supports coordinate descriptions of surfaces, equilibrium analysis, constrained optimization, and nonlinear equations.

\term{In-degree}
In a directed graph, the in-degree of a vertex $v$ is the number of edges pointing into $v$. In a regular tournament, all vertices have the same in-degree. The sum of in-degrees equals the sum of out-degrees, both equaling the total number of edges.

\term{Incenter}
The center of the incircle of a triangle, formed by the intersection of the internal angle bisectors. It is always equidistant from the three sides. The incenter is always located in the interior of the triangle.

\term{Incircle}
The largest circle that can be inscribed within a triangle, tangent to all three sides. Its radius $r$ is given by $r = \text{Area}/s$ where $s$ is the semi-perimeter.

\term{Independence}
Events $A$ and $B$ are independent when $P(A\cap B)=P(A)P(B)$, meaning learning one does not change the probability of the other. Random variables are independent when their joint distribution factors into the product of their marginal distributions. \par\noindent\textbf{Applications:} independence simplifies probability models, sampling, reliability calculations, and statistical inference.

\term{Inscribed Angle}
An angle formed by two chords that share an endpoint on a circle's circumference. An inscribed angle is always half the measure of the central angle subtending the same arc.

\term{Inscribed Circle}
Also known as the incircle. A circle that is tangent to all sides of a polygon. Only some polygons (like all triangles) possess an inscribed circle.

\term{Index (Group Theory)}
The index $[G:H]$ is the number of left or right cosets of a subgroup $H$ in $G$. If $G$ is finite, Lagrange's theorem gives $[G:H]=|G|/|H|$. \par\noindent\textbf{Applications:} index counts symmetry classes, orbit sizes, and the number of states obtained after identifying transformations in a subgroup.

\term{Index (Complex Analysis)}
The index, or winding number, of a closed curve $\gamma$ around $a\notin\gamma$ is
\[
\operatorname{ind}_\gamma(a)=\frac{1}{2\pi i}\oint_\gamma\frac{dz}{z-a}.
\]
It counts signed turns around the point. \par\noindent\textbf{Applications:} winding numbers support the argument principle, contour integration, and the study of zeros and poles of complex functions.

\term{Index (Fixed Point)}
A topological invariant assigned to a fixed point of a map $f: M \to M$. The sum of indices of all fixed points equals the Euler characteristic $\chi(M)$ by the Poincaré--Hopf Theorem.

\term{Indicator Function}
The indicator function $\mathbb 1_A(x)$ is $1$ when $x\in A$ and $0$ otherwise. Its integral records the measure of the set: $\int\mathbb1_A\,d\mu=\mu(A)$. \par\noindent\textbf{Applications:} indicators encode events, define piecewise functions, and simplify probability, measure theory, statistics, and optimization models.

\term{Inductive Proof}
A method of proof where a statement $P(n)$ is proved for a base case $n_0$, and then the implication $P(k) \implies P(k+1)$ is proved. This establishes $P(n)$ for all $n \geq n_0$.

\term{Induction (Mathematical)}
Mathematical induction proves a statement for every natural number by establishing a base case and showing $P(n)\Rightarrow P(n+1)$. Strong induction permits using all earlier cases $P(m)$ for $m<n$. \par\noindent\textbf{Applications:} induction proves formulas, recursive algorithms, divisibility results, and correctness of constructions defined step by step.

\term{Inductive Limit}
In category theory, the colimit of a directed system of objects. For a sequence of nested spaces $X_1 \subseteq X_2 \subseteq \cdots$, the inductive limit is their union $\bigcup X_i$. It is the categorical dual of the projective limit.

\term{Inertia Tensor}
The inertia tensor
\[
I_{ij}=\int\rho(r^2\delta_{ij}-r_ir_j)\,dV
\]
describes how mass is distributed relative to rotation axes. It relates angular momentum and angular velocity by $\mathbf L=I\boldsymbol\omega$; its eigenvectors are principal axes. \par\noindent\textbf{Applications:} it predicts rotation of rigid bodies, spacecraft attitude, robotics motion, and molecular dynamics.

\term{Inequation}
A statement asserting that two expressions are not equal, typically using $\leq$ or $\geq$. Solving an inequation involves finding the set of values for which the inequality holds.

\term{Inference}
The derivation of a conclusion from premises according to the rules of logic.

\term{Infimum}
The infimum $\inf S$ is the greatest number that is no larger than every element of $S$. It need not belong to $S$: the infimum of $(0,1)$ is $0$. \par\noindent\textbf{Applications:} infima define limits and optimization objectives when a minimum is not attained and are central to completeness and variational methods.

\term{Infinite Set}
A set that is not finite; it cannot be placed in bijection with any $\{1, \ldots, n\}$. A set is infinite iff it has a proper subset with the same cardinality (Dedekind-infinite).

\term{Infinite Series}
An infinite series $\sum_{n=1}^\infty a_n$ is defined through its partial sums $S_N=\sum_{n=1}^N a_n$. It converges when $S_N$ has a finite limit. The ratio, root, comparison, and integral tests help determine convergence. \par\noindent\textbf{Applications:} series approximate functions, solve differential equations, and quantify accumulated probabilities or errors.

\term{Infinite Product}
The limit of a sequence of partial products $\prod_{n=1}^\infty (1+a_n)$. It converges to a non-zero value iff $\sum a_n$ converges absolutely. Used in the Weierstrass product theorem.

\term{Infinite Dimensional Space}
A vector space is infinite-dimensional if no finite list of vectors spans it; examples include $L^p$ spaces and all polynomials. In such spaces, convergence and continuity require a chosen norm or topology. \par\noindent\textbf{Applications:} infinite-dimensional spaces model functions, signals, quantum states, and solutions of differential equations.

\term{Infinitesimal}
A quantity considered to be closer to zero than any standard real number, but not zero. In non-standard analysis, infinitesimals are formal elements of the hyperreal numbers ${}^*\mathbb{R}$.

\term{Infinitesimal Generator}
For a strongly continuous semigroup $T_t$ on a Banach space, the operator $A = \lim_{t \to 0} \frac{T_t - I}{t}$ defined on the domain where the limit exists. It represents the rate of change at $t=0$.

\term{Infinity}
A quantity larger than every finite number, denoted $\infty$, handled rigorously through limits and cardinalities.

\term{Information Entropy}
For a discrete random variable with probabilities $p_i$, Shannon entropy is $H(X)=-\sum_i p_i\log p_i$. It measures average uncertainty: rare outcomes carry more information, and entropy is largest for a uniform distribution with a fixed number of outcomes. \par\noindent\textbf{Applications:} entropy guides data compression, coding, decision trees, privacy, and statistical modeling.

\term{Information Theory}
The mathematical study of the quantification, storage, and communication of information, founded by Claude Shannon (1948). Key concepts include entropy, mutual information, and channel capacity.

\term{Information Gain}
Information gain is the decrease in entropy produced by observing or splitting on a variable. For a decision tree, the preferred split is often the one with the largest expected reduction in class uncertainty. \par\noindent\textbf{Applications:} it selects features in decision trees and helps quantify the usefulness of observations in data analysis.

\term{Information Geometry}
Information geometry treats a family of probability distributions as points on a manifold. The Fisher information matrix supplies a metric that measures how distinguishable nearby distributions are. \par\noindent\textbf{Applications:} it explains natural-gradient optimization, parameter estimation, exponential families, and statistical learning.

\term{Initial Condition}
Data specifying the state of a dynamical system at an initial time, used to select a unique solution of an ODE.

\term{Initial-value Problem}
An initial-value problem consists of a differential equation and data specifying the solution, and possibly its derivatives, at an initial point or time. Existence and uniqueness conditions determine whether those data select one solution. \par\noindent\textbf{Applications:} initial-value problems model motion, circuits, population change, heat flow, and control systems.

\term{Injective}
A function mapping distinct inputs to distinct outputs, $f(a) = f(b) \Rightarrow a = b$. Also called one-to-one.

\term{Injective Function}
A function $f:X\to Y$ is injective, or one-to-one, if $f(x_1)=f(x_2)$ implies $x_1=x_2$. Distinct inputs have distinct outputs, so $f$ has an inverse on its image even if it is not onto $Y$. \par\noindent\textbf{Applications:} injectivity guarantees recoverable encodings, unique parameter identification, and cancellation in algebra.

\term{Injective Module}
A module $M$ such that any embedding $M \subseteq N$ splits. Equivalently, the functor $\operatorname{Hom}(-, M)$ is exact. Every module can be embedded into an injective module.

\term{Injective Hull}
The smallest injective module containing a given module $M$, denoted $E(M)$. It is the unique (up to isomorphism) essential extension of $M$ that is injective.

\term{Injective Limit}
Another name for the inductive limit (colimit) of a directed system. It describes the smallest object that contains all objects in the system consistently.

\term{Inner Automorphism}
An inner automorphism of a group is conjugation by a fixed element: $g\mapsto aga^{-1}$. It preserves the group operation and describes a symmetry internal to the group. \par\noindent\textbf{Applications:} inner automorphisms organize conjugacy classes, normal subgroups, and the structure of symmetry groups.

\term{Inner Product Space}
An inner-product space is a vector space with a positive-definite form $\langle u,v\rangle$ that defines lengths $\|u\|=\sqrt{\langle u,u\rangle}$, angles, and orthogonality. In complex spaces the form is conjugate-linear in one argument. \par\noindent\textbf{Applications:} inner products support projections, Fourier expansions, least squares, and geometric models of data.

\term{Inner Product}
A mapping $\langle \cdot, \cdot \rangle: V \times V \to \mathbb{F}$ that is linear in the first argument, conjugate symmetric, and positive definite. It generalises the dot product in $\mathbb{R}^n$.

\term{Inseparable Extension}
An algebraic extension $L/K$ where some element has a minimal polynomial with multiple roots in its splitting field. This occurs only in characteristic $p > 0$.

\term{Inseparable Polynomial}
A polynomial that has multiple roots in its splitting field. For a field of characteristic $p$, a polynomial is inseparable iff it is a function of $x^p$.

\term{Inspection Paradox}
The inspection paradox is the tendency for a randomly encountered interval to be longer than a typical interval. For example, arriving at a random time is more likely to occur during a long bus gap than a short one. It is caused by length-biased sampling. \par\noindent\textbf{Applications:} it matters in queueing, reliability, renewal processes, traffic, and service-wait analysis.

\term{Instanton}
A solution of the Yang--Mills equations on a principal $G$-bundle over a Riemannian four-manifold with finite action, equivalently a connection whose curvature $F$ satisfies the self-duality equation $F = *F$ (or anti-self-duality $F = -*F$), where $*$ is the Hodge star. Instantons minimise the Yang--Mills functional in their topological class, and their classification by Chern classes and Pontryagin number uses index theory; over $\mathbb{R}^4$ the ADHM construction parametrises all instantons. Donaldson invariants of smooth four-manifolds are built from moduli spaces of instantons, making them central to gauge theory and the differential topology of four-manifolds.

\term{Integers}
The integers are $\mathbb Z=\{\ldots,-2,-1,0,1,2,\ldots\}$. They form a commutative ring and Euclidean domain, so division with remainder and unique prime factorization are available. The only multiplicative units are $\pm1$. \par\noindent\textbf{Applications:} integers model discrete quantities and form the basis of number theory, modular arithmetic, coding, and exact computation.

\term{Integer Partition}
An integer partition writes a positive integer $n$ as a sum of positive integers, with order ignored; for example, $4=3+1=2+2=2+1+1=1+1+1+1$ gives four nontrivial forms plus $4$ itself. The partition number $p(n)$ grows rapidly. \par\noindent\textbf{Applications:} partitions count combinations in number theory, statistical mechanics, representation theory, and algorithmic enumeration.

\term{Integer Programming}
Integer programming optimizes a linear or nonlinear objective subject to constraints while requiring some or all variables to be integers. Integrality makes many problems combinatorial and NP-hard. \par\noindent\textbf{Applications:} it models scheduling, facility location, routing, production planning, portfolio selection, and network design.

\term{Integral}
An integral assigns a limit to accumulated quantities. The Riemann integral uses partitions of the domain, while the Lebesgue integral groups points by function values and handles more general functions. \par\noindent\textbf{Applications:} integrals compute area, mass, probability, energy, work, and expected values.

\term{Integral Calculus}
The branch of calculus dealing with the computation of integrals and the area under curves, linked to differential calculus via the Fundamental Theorem of Calculus.

\term{Integral Equation}
An integral equation contains an unknown function under an integral sign. Fredholm equations have fixed limits, while Volterra equations have a variable limit such as $\int_a^x$. Kernels describe how values at different points interact. \par\noindent\textbf{Applications:} integral equations model boundary-value problems, inverse problems, scattering, population dynamics, and transport.

\term{Integral Curve}
An integral curve of a vector field $V$ satisfies $\dot\gamma(t)=V(\gamma(t))$, so its tangent agrees with the field at every point. It is the trajectory of the corresponding first-order ODE. \par\noindent\textbf{Applications:} integral curves describe particle motion, fluid streamlines, dynamical systems, and control trajectories.

\term{Integral Manifold}
A submanifold whose tangent space at every point is the subspace specified by a distribution. By Frobenius' Theorem, a distribution is integrable iff it is involutive.

\term{Integrability}
In analysis, a function is integrable when its absolute integral is finite, for example $\int|f|<\infty$ in the Lebesgue sense. In dynamics, a system is integrable when enough independent conserved quantities reduce its motion to a solvable form. \par\noindent\textbf{Applications:} integrability controls whether totals are finite and whether differential equations admit explicit or structured solutions.

\term{Integrable Function}
A function $f$ for which the integral $\int f$ exists and is finite. In the Lebesgue sense, $f$ is integrable iff $\int |f| \, d\mu < \infty$.

\term{Integrable System}
A dynamical system with as many independent constants of motion as degrees of freedom. Their trajectories are confined to invariant tori (Liouville-Arnold theorem).

\term{Integrating Factor}
An integrating factor is a nonzero function $\mu$ that turns a differential equation into an exact one. For $y'+P(x)y=Q(x)$, the factor $\mu(x)=e^{\int P(x)dx}$ gives $(\mu y)'=\mu Q$. \par\noindent\textbf{Applications:} integrating factors solve first-order ODEs and appear in conservation laws and differential-form calculations.

\term{Integration by Parts}
Integration by parts is the integral form of the product rule:
\[
\int u\,dv=uv-\int v\,du.
\]
It transfers a derivative from one factor to another. \par\noindent\textbf{Applications:} it evaluates products of functions, derives weak PDE formulations, proves orthogonality relations, and establishes estimates.

\term{Integration by Substitution}
A method for simplifying integrals by changing variables, based on the chain rule. It transforms the integral $\int f(g(x))g'(x)\,dx$ into $\int f(u)\,du$.

\term{Integral Operator}
An operator $T$ defined by $(Tf)(x) = \int_a^b K(x,y) f(y)\,dy$. The function $K(x,y)$ is the kernel. Many such operators are compact on $L^2$.

\term{Integral Transform}
An integral transform maps a function to another function through an integral, as in the Fourier, Laplace, or Mellin transform. A suitable transform can turn differentiation or convolution into multiplication. \par\noindent\textbf{Applications:} transforms solve differential equations, analyze signals, model systems, and study asymptotic behavior.

\term{Integral Domain}
An integral domain is a commutative ring with $1\ne0$ and no zero divisors: $ab=0$ implies $a=0$ or $b=0$. It supports cancellation and embeds in a field of fractions. \par\noindent\textbf{Applications:} integral domains provide the algebraic setting for polynomial factorization, divisibility, and number theory.

\term{Integral Basis}
A basis for the ring of integers $\mathcal{O}_K$ of a number field $K$ as a $\mathbb{Z}$-module. If $\{\omega_i\}$ is an integral basis, every $\alpha \in \mathcal{O}_K$ is $\sum a_i \omega_i$ with $a_i \in \mathbb{Z}$.

\term{Integral Closure}
For a subring $R \subseteq S$, the set of elements in $S$ that are roots of monic polynomials in $R[x]$. $S$ is integrally closed if $S$ equals its closure in its fraction field.

\term{Integrally Closed Domain}
An integral domain that is equal to its integral closure in its field of fractions. Examples include UFDs and Dedekind domains.

\term{Integer-valued Polynomial}
A polynomial $p(x) \in \mathbb{Q}[x]$ such that $p(n) \in \mathbb{Z}$ for all $n \in \mathbb{Z}$. The binomial coefficients $\binom{x}{k}$ form a basis for such polynomials.

\term{Interior (Topology)}
The interior of a set $S$, denoted $\operatorname{int}(S)$, is the union of all open sets contained in $S$. It consists of all interior points of $S$.

\term{Interior Point}
A point $x \in S$ such that there exists an open neighborhood $U$ with $x \in U \subseteq S$.

\term{Interior Product}
An operation that contracts a vector field $X$ with a differential $k$-form $\omega$, resulting in a $(k-1)$-form: $\iota_X \omega$. It is a derivation on the exterior algebra.

\term{Interior Point Method}
Interior-point methods solve optimization problems by staying inside the feasible region and approaching its boundary through barrier functions. For linear programming they can achieve polynomial-time complexity under standard models. \par\noindent\textbf{Applications:} they solve large-scale resource allocation, network, engineering, and machine-learning optimization problems.

\term{Intermediate Value Theorem}
If $f$ is continuous on $[a,b]$ and $y$ lies between $f(a)$ and $f(b)$, then some $c\in[a,b]$ satisfies $f(c)=y$. A sign change therefore guarantees a zero. \par\noindent\textbf{Applications:} the theorem justifies bisection methods, existence of equilibria, and root-finding in scientific computation.

\term{Intercept (Curve)}
The points where a curve intersects the coordinate axes. For a function $y=f(x)$, the $y$-intercept is $f(0)$ and the $x$-intercepts are the roots of $f(x)=0$.

\term{Interpolation}
Interpolation constructs a function that matches given data points, using lines, polynomials, splines, or other bases. It estimates values between measured samples, unlike extrapolation outside the data range. \par\noindent\textbf{Applications:} interpolation supports numerical tables, computer graphics, signal reconstruction, and scientific data processing.

\term{Interpolating Polynomial}
The unique polynomial of degree at most $n-1$ that passes through $n$ given points. It can be constructed using Lagrange or Newton forms.

\term{Interpolation Error}
The difference between the actual function and its interpolant. For polynomials, it is bounded by a term involving the $(n+1)$-th derivative of the function.

\term{Interpolation Formula}
A general expression for an interpolant, such as the Lagrange formula $\sum y_i L_i(x)$ or the Newton divided-difference formula.

\term{Interpolation Space}
Spaces arising in interpolation theory, the study of how a linear operator bounded on two Banach spaces behaves on intermediate spaces. The real method, built on the $K$-functional, and the complex method construct the intermediate spaces $(A_0, A_1)_{\theta,q}$ and $[A_0, A_1]_\theta$ between a pair of Banach spaces. Interpolation inequalities such as the Riesz--Thorin theorem bound the norm of an operator on the interpolant in terms of its endpoint norms. Sobolev spaces and $L^p$ spaces arise as real interpolation spaces, and the theory is essential in harmonic analysis and the regularity theory of PDEs.

\term{Interquartile Range}
The interquartile range is $IQR=Q_3-Q_1$, the difference between the 75th and 25th percentiles. It measures the spread of the middle half of a dataset and is resistant to extreme outliers. \par\noindent\textbf{Applications:} it summarizes skewed data and is used to flag potential outliers with box plots and robust statistics.

\term{Intersection Theory}
The study of the intersection of subvarieties and cycles on a scheme, whose equivalence classes form the Chow ring, graded by codimension, with the intersection product corresponding to proper intersection of cycles. Rational equivalence, Chow groups, and the pushforward and pullback of cycles are the basic structures. Chern classes of vector bundles are computed via intersections with zero loci, and the theory governs enumerative geometry and birational geometry, connecting to algebraic K-theory through the relation between Chow groups and K-cohomology.

\term{Interval (Analysis)}
A connected subset of $\mathbb{R}$. Closed intervals $[a,b]$ are compact; open intervals $(a,b)$ are open. Intervals are the building blocks of the Borel $\sigma$-algebra.

\term{Interval Arithmetic}
A method for computing with intervals $[a,b]$ rather than point values, used to bound rounding errors and prove inequalities rigorously.

\term{Inverse Function}
A function $f^{-1}$ such that $f^{-1}(f(x)) = x$ and $f(f^{-1}(y)) = y$. $f$ must be bijective for $f^{-1}$ to exist on the whole codomain.

\term{Inverse Map}
A synonym for the inverse function in category theory or topology. A map $f: X \to Y$ is an isomorphism iff it has an inverse map $f^{-1}: Y \to X$.

\term{Inverse Element}
In a group, the element $g^{-1}$ such that $g \cdot g^{-1} = e$. In a ring, the multiplicative inverse $u^{-1}$ such that $u u^{-1} = 1$.

\term{Inverse Operator}
An operator $T^{-1}$ such that $T T^{-1} = T^{-1} T = I$. For matrices, this is the inverse matrix. In functional analysis, the inverse is often a bounded operator.

\term{Inverse Matrix}
A square matrix $A^{-1}$ such that $A A^{-1} = I$. It exists iff $\det(A) \neq 0$. The inverse is computed via the adjugate matrix or Gaussian elimination.

\term{Inverse Transform}
An operation that recovers the original function from its transform, e.g., the Inverse Fourier Transform $f(x) = \int \hat{f}(\xi) e^{2\pi i x \xi} \, d\xi$.

\term{Inverse Problem}
A problem where the goal is to determine the cause (e.g., a source) from the observed effects (e.g., the field). Often ill-posed and requires regularisation.

\term{Inverse Limit}
The categorical limit of a projective system of objects. It is the dual of the inductive limit. Examples include the $p$-adic integers $\mathbb{Z}_p = \varprojlim \mathbb{Z}/p^n\mathbb{Z}$.

\term{Invariant (Algebra)}
An expression or property that remains unchanged under a group action. For example, the determinant of a matrix is invariant under basis change.

\term{Invariant (Geometry)}
A property of a geometric object that is preserved under a group of transformations, such as distance under isometry or angle under conformal mapping.

\term{Invariant (Topology)}
A property of a space preserved under homeomorphism, such as the Euler characteristic or fundamental group. If two spaces have different invariants, they are not homeomorphic.

\term{Invariant Subspace}
A subspace $W \subseteq V$ such that $T(W) \subseteq W$ for a linear operator $T$. The goal of many decomposition theorems (e.g., Jordan form) is to find a basis of invariant subspaces.

\term{Invariant Measure}
A measure $\mu$ on a space $X$ that is preserved by a transformation $T$: $\mu(T^{-1} A) = \mu(A)$ for all measurable $A$. Ergodic theory studies these measures.

\term{Invariant Set}
A set $S$ in a dynamical system such that if $x \in S$, then $f(x) \in S$ for all time. Examples include attractors and unstable manifolds.

\term{Invariant Ring}
The ring of all polynomials $p \in K[x_1, \ldots, x_n]$ that are invariant under the action of a group $G$. The study of such rings is the core of invariant theory.

\term{Invariant Theory}
The branch of algebra studying polynomial invariants. Classical invariant theory focuses on SL(n) and GL(n) acting on tensors.

\term{Inverse Square Law}
A physical law stating that a specified physical quantity is inversely proportional to the square of the distance from the source, e.g., Newton's law of gravity.

\term{Involutive Distribution}
A distribution $D$ such that for any two vector fields $X, Y \in D$, their Lie bracket $[X, Y]$ is also in $D$. By Frobenius' Theorem, such distributions are integrable.

\term{Involution (Group Theory)}
An element $g$ of a group such that $g^2 = e$ (and $g \neq e$). The set of involutions often plays a key role in the classification of finite simple groups.

\term{Involution (Algebra)}
An antiautomorphism $\ast: A \to A$ of an algebra such that $(a^\ast)^\ast = a$ and $(ab)^\ast = b^\ast a^\ast$. C*-algebras are defined via such involutions.

\term{Involutory Matrix}
A square matrix $A$ such that $A^2 = I$. Such matrices represent reflections in geometry and have eigenvalues $\pm 1$.

\term{Icosahedron}
A regular polyhedron with 20 equilateral triangular faces, 12 vertices, and 30 edges. It is one of the five Platonic solids.

\term{Icosahedral Group}
The symmetry group of the regular icosahedron, which is isomorphic to the alternating group $A_5$ (for rotational symmetries) or $A_5 \times \mathbb{Z}/2\mathbb{Z}$.

\term{Irrational Number}
A real number that is not rational, i.e. not expressible as a ratio of integers, such as $\sqrt{2}$ or $\pi$.

\term{Irreducible}
An element of a ring that cannot be written as a product of two non-units; an irreducible polynomial cannot be factored over the ground field.

\term{Isolated Point}
A point of a set having a neighbourhood containing no other points of the set.

\term{Isomorphism}
A bijective map $f: X \to Y$ that preserves the structure of the objects. In groups, $f(ab) = f(a)f(b)$. In topology, it is a homeomorphism.

\term{Isometry}
A distance-preserving map between metric spaces: $d(f(x), f(y)) = d(x,y)$. Isometries are always injective. Examples include translations, rotations, and reflections.

\term{Isoperimetric Inequality}
The inequality $4\pi A \leq L^2$ relating the area $A$ of a plane region to its perimeter $L$, with equality for circles.

\term{Isothermal Coordinates}
Local coordinates in which a surface's metric takes the conformal form $\lambda^2(dx^2 + dy^2)$.

\term{Iteration}
The repeated application of a function or process; in numerical analysis, the basis of fixed-point methods.

\term{Iterative Method}
A numerical method that generates successive approximations converging to the solution.

\term{Itô Calculus}
The stochastic calculus for semimartingales, centered on the Itô integral of predictable processes with respect to martingales such as Brownian motion. Itô's formula is the change-of-variables rule, whose second-order term is the quadratic variation $[X,X]_t$, and it is the source of stochastic differential equations, whose solutions are diffusions governed by their generators. Unlike the Stratonovich integral, the Itô integral is an isometry and a martingale. Itô calculus is the foundation of the theory of diffusions and of mathematical finance through the Itô isometry and Itô's lemma.
\term{Finite Element Method}
A numerical technique for solving PDEs by partitioning the domain into simple elements and approximating the solution on each element by basis functions, typically polynomials. The method converts a PDE into a large sparse linear or nonlinear system. It is the dominant method in engineering simulation and has a rigorous foundation in Sobolev space theory and Céa's lemma.
\term{Initial Object}
An object $0$ in a category such that for every object $X$ there is a unique morphism $0 \to X$. Dually, a terminal object has a unique morphism from every object. In **Set**, the initial object is the empty set and the terminal objects are singleton sets. In **Grp**, the trivial group is both initial and terminal.
\term{Independent Set}
A set of vertices in a graph no two of which are adjacent. The independence number $\alpha(G)$ is the maximum size of an independent set. Independent sets are dual to vertex covers: a set is a vertex cover if and only if its complement is an independent set. The problem of finding maximum independent sets is NP-hard.
